\section*{Projectile Simulation}

There is a projectile which is launched with
\[
\vec{r_0} = (x_0, y_0)
\]
\[
\vec{v_0} = (v_{0x}, v_{0y})
\]

The projectile is affected by gravity:
\[
\vec{F_g} = (0, -mg)
\]

The projectile is also affected by air resistance:
\[
\vec{F_r} = (-k v_x, -k v_y)
\]

There are some lines that the projectile can collide with. Each line is represented by two points:
\[
p_i = (\vec{p_{i1}}, \vec{p_{i2}})
\]

And there is constraint that
\[
\forall i, \vec{d_i} \parallel \hat{x} \lor \vec{d_i} \parallel \hat{y}
\]
\[
\forall i, \vec{p}_{i2} = \vec{p}_{(i+1)1}
\]

where $\vec{d_i} = \vec{p_{i2}} - \vec{p_{i1}}$.


\vspace{3em}
\hrule
\vspace{2em}

\section*{Numerical Integration Convergence}

In physics simulations, we often use simple update rules to approximate continuous motion. This section explains that as the time step $dt \to 0$, these numerical updates converge to the true analytical physical solutions.

\subsection*{Case 1: Drag Force (Velocity Convergence)}

\textbf{Mathematical Model:}
Consider a particle subject to a drag force proportional to its velocity:
\[
F = -k v \implies m \frac{dv}{dt} = -k v \implies \frac{dv}{dt} = -\frac{k}{m} v
\]
The analytical solution is $v(t) = v_0 e^{-(k/m)t}$.

\textbf{Discretized Update Rule:}
In a simulation, we update the velocity using the rule:
\[
v += -(k/m) v \cdot dt \implies v_{n+1} = v_n \left(1 - \frac{k}{m} dt\right)
\]

\textbf{Convergence Explanation:}
After $n$ steps ($t = n \cdot dt$), the velocity is:
\[
v_n = v_0 \left(1 - \frac{k}{m} dt\right)^n = v_0 \left(1 - \frac{k}{m} \frac{t}{n}\right)^n
\]
Using the definition of the exponential function $\lim_{n \to \infty} (1 - \frac{a}{n})^n = e^{-a}$:
\[
\lim_{dt \to 0} v_n = v_0 e^{-(k/m)t}
\]
Thus, the simulation result converges to the actual physical solution as $dt \to 0$.

\subsection*{Case 2: Integrating Velocity to Position}

\textbf{Mathematical Model:}
Given a velocity function $v(t)$, the continuous equation is $dx/dt = v(t)$, which yields:
\[
x(t) = x_0 + \int_0^t v(s) ds
\]

\textbf{Discretized Update Rule:}
The simulation repeatedly applies:
\[
x += v \cdot dt \implies x_{n+1} = x_n + v(t_n) dt
\]

\textbf{Convergence Explanation:}
Summing the updates yields:
\[
x_n = x_0 + \sum_{k=0}^{n-1} v(t_k) dt
\]
This is exactly the \textbf{Riemann Sum} of the velocity function. As $dt \to 0$, the Riemann sum converges to the definite integral:
\[
\lim_{dt \to 0} x_n = x_0 + \int_0^t v(s) ds = x(t)
\]

\subsection*{Case 3: Full Trajectory (Acceleration to Position)}

\textbf{Mathematical Model:}
Given an acceleration $a(t)$, the system is:
\[
\frac{dv}{dt} = a(t), \quad \frac{dx}{dt} = v(t)
\]

\textbf{Discretized Update Rule:}
The simulation updates both velocity and position:
\[
v += a \cdot dt, \quad x += v \cdot dt
\]

\textbf{Convergence Explanation:}
As shown in Case 2, the velocity update converges to $v(t) = v_0 + \int a(t) dt$. Then, the position update, which uses these converged velocity values, also converges to the integral of $v(t)$. Therefore, the entire chain $a \to v \to x$ converges to the physical trajectory as $dt \to 0$.
